\DocumentMetadata{
  pdfversion=2.0,
  pdfstandard=ua-2,
  lang=en-US,
  tagging=on,
  tagging-setup={math/setup=mathml-SE,tabsorder=structure}
}
\documentclass[11pt,letterpaper]{article}
\usepackage[margin=0.68in,includefoot]{geometry}
\usepackage{newcomputermodern}
\usepackage{amsmath,amscd,mathtools}
\providecommand{\square}{\mdlgwhtsquare}
\usepackage[hidelinks]{hyperref}
\hypersetup{
  pdftitle={Ergodic Theory PhD Exam, May 2013},
  pdfauthor={University of Florida Department of Mathematics},
  pdfsubject={Graduate mathematics examination},
  pdfdisplaydoctitle=true,
  bookmarksopen=true
}
\setcounter{secnumdepth}{0}
\setlength{\parindent}{0pt}
\setlength{\parskip}{0.28em}
\setlength{\emergencystretch}{3em}
\setlength{\leftmargini}{2.15em}
\setlength{\leftmarginii}{2.65em}
\setlength{\leftmarginiii}{3em}
\renewcommand{\labelenumi}{\textbf{\theenumi.}}
\pagestyle{empty}
\raggedbottom
\allowdisplaybreaks
\newenvironment{examproblems}[1]{%
  \begin{enumerate}%
  \setcounter{enumi}{#1}%
  \setlength{\topsep}{0.35em}%
  \setlength{\partopsep}{0pt}%
  \setlength{\itemsep}{0.48em}%
  \setlength{\parsep}{0.18em}%
}{\end{enumerate}}
\newenvironment{examsubparts}{%
  \begin{enumerate}%
  \setlength{\topsep}{0.18em}%
  \setlength{\partopsep}{0pt}%
  \setlength{\itemsep}{0.16em}%
  \setlength{\parsep}{0.1em}%
}{\end{enumerate}}
\newenvironment{examinstructions}{%
  \par\begingroup\itshape
}{%
  \par\endgroup\vspace{0.18em}
}
\makeatletter
\renewcommand\section{\@startsection{section}{1}{0pt}%
  {-1.25ex plus -.25ex minus -.1ex}%
  {0.55ex plus .1ex}%
  {\normalfont\large\bfseries}}
\renewcommand{\@maketitle}{%
  \newpage\null\vskip -1.2em%
  {\footnotesize\bfseries
    \href{https://www.ufl.edu/}{University of Florida}, \href{https://math.ufl.edu/}{Department of Mathematics}\par}%
  \vskip 0.18em%
  \tagstructbegin{tag=Title}%
    {\LARGE\bfseries\href{https://gma.math.ufl.edu/exams/ergodic-theory-phd/}{Ergodic Theory PhD Exam}, May 2013\par}%
  \tagstructend%
  \vskip 0.35em%
  \tagmcbegin{artifact}%
    \hrule height 1pt%
  \tagmcend%
  \vskip 0.55em%
}
\makeatother
\title{Ergodic Theory PhD Exam, May 2013}
\author{}
\date{}
\begin{document}
\maketitle
\thispagestyle{empty}
\begin{examproblems}{0}
\item
Suppose that~\(X\) is a compact metric space, and \(f:X\to X\) is continuous. Suppose that~\(X\) is a minimal set for~\(f\). Prove that every point of~\(X\) is almost periodic.
\item
Let \(f:[0,1]\to[0,1]\) be continuous. Suppose that~\(f\) has a periodic point of odd period larger than one. Prove that the topological entropy of~\(f\) is at least \(\frac{\log(2)}{2}\).
\item
Describe the construction of the homeomorphism of the sphere to itself known as the Horseshoe Map (of S. Smale). Prove that the set of periodic points is dense in the nonwandering set for this map.
\item
Let~\(X\) be a topological space, and let \(f:X\to X\) be a Borel measurable function. Suppose that there exists a unique~\(f\)-invariant Borel probability measure, \(\mu\). Prove that~\(\mu\) is ergodic.
\item
Let \(T_1:[0,1]\to[0,1]\) denote the full tent map, let \(\mathcal B_1\) denote the collection of Borel sets, and let \(\mu_1\) denote Lebesgue measure. Recall that \(T_1\) is measure preserving with respect to the probability space \(([0,1],\mathcal B_1,\mu_1)\). Let \(T_2:\Sigma_2^+\to\Sigma_2^+\) denote the full one-sided shift map on two symbols, let \(\mathcal B_2\) denote the~\(\sigma\)-algebra generated by the cylinder sets, and let \(\mu_2\) denote the \((\frac12,\frac12)\) product measure. Recall that \(T_2\) is measure preserving with respect to the probability space \((\Sigma_2^+,\mathcal B_2,\mu_2)\). Are \(T_1\) and \(T_2\) isomorphic? Prove your answer.
\item
Suppose that~\(X\) is a compact metric space, and \(f:X\to X\) is continuous. Prove that there exists an~\(f\)-invariant Borel probability measure.
\item
Let \((X,\mathcal B,\mu)\) be a probability space. Suppose that \(T:X\to X\) is measure-preserving. Prove that the following are equivalent.
\begin{examsubparts}
\item[\textnormal{(a)}]
\(T\) is ergodic.
\item[\textnormal{(b)}]
For every~\(f\) in \(L^1(\mu)\), the sequence of real numbers 
\[
\left\|\int f\,d\mu-\frac1n\sum_{k=0}^{n-1}f\circ T^k\right\|_1
\]
 approaches~\(0\) as \(n\to\infty\).
\item[\textnormal{(c)}]
For every bounded~\(f\) in \(L^1(\mu)\), the sequence of real numbers 
\[
\left\|\int f\,d\mu-\frac1n\sum_{k=0}^{n-1}f\circ T^k\right\|_1
\]
 approaches~\(0\) as \(n\to\infty\).
\end{examsubparts}
\item
Let \(f:[0,2]\to[0,2]\) be the piecewise linear continuous map given by \(f(x)=x+1\), if \(0\le x\le1\) and \(f(x)=4-2x\), if \(1\le x\le2\). Determine whether each of the following statements is true or false. Prove your conclusions.
\begin{examsubparts}
\item[\textnormal{(a)}]
Lebesgue measure is an invariant measure for~\(f\).
\item[\textnormal{(b)}]
There exist infinitely many (distinct) invariant ergodic measures for~\(f\).
\item[\textnormal{(c)}]
The map~\(f\) has a periodic point with period \(837\).
\item[\textnormal{(d)}]
The map~\(f\) is forward orbit topologically transitive.
\item[\textnormal{(e)}]
Each point of \([0,2]\) is nonwandering under~\(f\).
\end{examsubparts}
\end{examproblems}
\end{document}
